\section{関連研究}

Dongら\cite{dong2019}は，NVDの標準化されたエントリと，CVE概要文および参照先の非構造な脆弱性レポートとの不整合を検出するシステムVIEMを提案した．
深層学習に基づく固有表現抽出と関係抽出により製品名とバージョンを抽出し（precision 0.941, recall 0.993），約20年分・78,296件のCVE IDと70,569件のレポートを対象とした調査で，NVDエントリと影響バージョンが完全に一致するのは59.82\%にすぎないことを示した．

本研究はDongらと問題意識を共有するが，3点で異なる．
第一に，VIEMはNVDにエントリが存在することを前提に不整合を測るのに対し，本稿はCVE増加後に顕在化した\textbf{検知条件がそもそも生成されない}という失敗モード（\texttt{Deferred} / \texttt{Awaiting Analysis}）を対象とする．
第二に，VIEMは専用モデルの学習と表記変化ごとの再学習を要するのに対し，本稿は抽出精度単体ではなく\textbf{継続運用における保守性・再現性}を設計目標に置く．
第三に，VIEMは検出までを扱うのに対し，本稿は不整合を解消したデータベースを構築し検知精度への影響まで評価する．

NVDのデータ品質については，Anwarら\cite{anwar2022}が公開日・ベンダ名・深刻度等の不整合と欠損を体系的に調査し，NguyenとMassacci\cite{nguyen2013}がGoogle Chromeを対象に影響バージョンデータの系統的な誤りを示している．
いずれもCVE増加以前の調査であり，本稿は同種の問題が2024年以降に「欠損」として急拡大したことを示す点で相補的である．
